\documentclass{article}
\usepackage{graphicx} % Required for inserting images
\usepackage{float}

\title{Sistemi Operativi}
\author{Tommaso Trentin}
\date{June 2025}

\begin{document}

\newpage
\tableofcontents
\newpage
\section{Generalità}
\subsection{Gestione}
Problematiche da affrontare:
\begin{itemize}
    \item \textbf{allocazione} singoli job / processi in memoria
    \item \textbf{protezione} spazio di indirizzamento
    \item \textbf{condivisione} spazio di indirizzamento
    \item gestione \textbf{swapping}
\end{itemize}
\textbf{Vincolo di funzionamento}: per essere eseguito un programma deve essere portato in memoria e \textit{trasformato} in un processo:
\begin{itemize}
    \item CPU preleva istruzioni da eseguire da memoria in base a valore del P.C.
    \item ogni istruzione decodificata, può prevedere prelievo di operandi dalla memoria
    \item al termine di esecuzione dell’istruzione risultato può dover essere scritto in memoria
    \item al termine del processo la “sua” memoria viene rilasciata
\end{itemize}
\subsection{Indirizzi di memoria}
Trasformazione programma in processo implica assunzioni di diversa semantica di indirizzi e spazio di memoria $\rightarrow$ introduzione distinzione tra \textbf{spazio indirizzi logici} e \textbf{spazio indirizzi fisici}.
\subsubsection{Binding degli indirizzi}
\begin{itemize}
    \item indirizzi del programma sorgente sono simbolici
    \item durante compilazione si associano a indirizzi simbolici indirizzi \textbf{rilocabili}
    \item in fase collegamento o caricamento si associano a indirizzi rilocabili indirizzi assoluti
\end{itemize}
\includegraphics[width=0.6\textwidth]{OS_images/BInding_Indirizzi.png} \\
\textbf{Il binding di dati/istruzioni a indirizzi assoluti può avvenire in diversi momenti:} 
\begin{itemize}
    \item Al tempo di compilazione (\textit{binding statico al compile time}):
        \begin{itemize}
            \item Se si conosce a priori la posizione in memoria, si può generare un \textbf{codice assoluto}.
            \item Se si vuole cambiare posizione, è necessario \textbf{ricompilare}.
        \end{itemize}
    \item Al tempo di caricamento (\textit{binding statico al load time}):
        \begin{itemize}
            \item Serve generare un \textbf{codice rilocabile} (ripositionabile).
            \item Gli indirizzi sono \textbf{relativi} all’inizio del programma.
            \item Per cambiare posizione, basta \textbf{ricaricare}.
        \end{itemize}
    \item Al tempo di esecuzione (\textit{binding dinamico al run time}):
        \begin{itemize}
            \item Il binding è \textbf{posticipato}, il processo può essere spostato anche durante l’esecuzione.
            \item Richiede \textbf{supporto hardware} (es. registri base e limite).
        \end{itemize}
\end{itemize}
\includegraphics[width=0.6\textwidth]{OS_images/Binding_Indirizzi2.png} \\
Un HW particolare effettua il binding dinamico: la \textbf{MMU}
\subsubsection{MMU}
Contenuto di un particolare registro \textbf{registro di rilocazione} viene aggiunto a ogni indirizzo riferito da un processo, con l'indirizzo risultante si accede alla memoria: \\
\includegraphics[width:0.6\textwidth]{OS_images/Semplice_MMU.png}
\subsubsection{Caricamento}
\begin{itemize}
    \item Statico: intero codice caricato in memoria per procedere all'esecuzione
    \item Dinamico: cerca di migliorare utilizzo della memoria centrale
        \begin{itemize}
            \item moduli di codice caricati al momento del primo utilizzo
            \item conseguenza: parti di codice non utilizzate non vengono caricate $\rightarrow$ "risparmio" memoria
            \item tecnica utile con codici suddivisibili in parti che gestiscono casi specifici (error handling)
        \end{itemize}
\end{itemize}
Solitamente caricamento dinamico richiede intervento minimo da parte del S.O. (può essere realizzato a livello di linguaggio/programma)
\subsubsection{Collegamento}
\begin{itemize}
    \item Statico: tutti i riferimenti del codice definiti prima del'avvio dell'esecuzione. Immagine del processo contiene interamente il codice di eventuali librerie
    \item Dinamico: utilizzato solamente per il codice delle librerie
        \begin{itemize}
            \item linking posticipato al tempo di esecuzione
            \item conseguenza: codice del programma non include codice delle librerie, ma solo dei riferimenti (\textit{stub})
            \item stub indica dove trovare codice della libreria alla prima necessità (una volta caricata stub viene rimpiazzato con sua locazione in memoria)
            \item tutti i processi eseguono un'istanza dello stesso codice di librerie (singola copia in memoria) 
        \end{itemize}
\end{itemize}
Linking dinamico richiede intervento essenziale del S.O. (può interessare spazi di memoria di processi diversi).
\subsubsection{Swapping}
Tecnica che permette di rimuovere un processo da memoria principale e spostarlo in memoria secondaria (\textbf{swapped out}), in seguitoquesto potrà essere ricaricato in memoria principale (\textbf{swapped in}). \\
\includegraphics[width=0.6\textwidth]{OS_images/Swapping.png}
\begin{itemize}
    \item tempo di swapping proporzionale a ammontare dei dati trasferiti
    \item sensato in caso di multiprogrammazione (altrimenti dovrebbe essere attivato ad ogni context switch)
    \item di norma è necessario supporto a binding dinamico per poter ri-locare i processi al momento dello swap-in
    \item con binding statico invece va ricaricato il processo sempre nella stessa posizione
    \item moderni S.O. contemplano lo swapping (in qualche forma)
\end{itemize}
La tecnica definita finora è detta \textbf{swapping standard}, mentre alcuni S.O. ne adottano forme diverse:
\begin{itemize}
    \item per evitare di trasferire intero processo molti S.O. moderni effettuano swapping solo su porzioni del processo (\textit{pagine}) $\rightarrow$ tecnica connessa a uso della \textit{paginazione}
    \item S.O. di dispositivi mobili non operano swapping di processi, poiché:
        \begin{itemize}
            \item memoria di solito è flash, assenza di HD
            \item memoria flash presenta un limite al numero di scritture, swapping degraderebbe la memoria
        \end{itemize}
    \item se si necessita liberare memoria il S.O., invece di effettuare swapping, richiede a app di rinunciare a parte di memoria allocata (o sottrae a app pagine di memoria non allocate / termina la app)
\end{itemize}
\subsubsection{Considerazioni}
\begin{itemize}
    \item in presenza di multiprogrammazione è impossibile fissare staticamente locazione ove processi vengono caricati
    \item presenza di swapping impone impiego di rilocazione dinamica
\end{itemize}
Sono necessarie tecniche specifiche per gestione memoria principale, in caso il programma sia \textbf{interamente} caricato in memoria principale le tecniche "main" sono:
\begin{itemize}
    \item Allocazione contigua
    \item Paginazione 
    \item Segmentazione
    \item Segmentazione Paginata
\end{itemize}
\section{Allocazione Contigua}
\subsection{Concetti Generali}
\begin{itemize}
    \item processi allocati in memoria in posizioni \textbf{contigue} (ogni processo occupa un'\textbf{unica} sezione di memoria pari alla sua dimensione)
    \item la memoria viene suddivisa in due "zone":
        \begin{itemize}
            \item una utilizzata dal S.O. (contenente vettore delle interruzioni)
            \item una dedicata ai processi utente
        \end{itemize}
    \item memoria quindi suddivisa in partizioni disgiunte (ognuna accoglierrà un processo), due varianti:
        \begin{itemize}
            \item Allocazione contigua con \textbf{partizioni fisse}
            \item Allocazione contigua con \textbf{partizioni variabili}
        \end{itemize}
\end{itemize}
\subsection{Partizioni Fisse}
\begin{minipage}{0.45\textwidth}
\begin{itemize}
    \item Memoria suddivisa in un insieme di partizioni di dimensioni prefissate (non necessariamente uguali)
    \item Scelta di quante partizioni / dimensioni viene fatta una volta per tutte
    \item Un processo viene caricato in una partizione in grado di contenerlo interamente
\end{itemize}
\end{minipage}
\hfill
\begin{minipage}{0.5\textwidth}
    \includegraphics[width=\textwidth]{OS_images/Partizioni_Fisse.png}
\end{minipage}
\subsubsection{Allocazione della Memoria}
Long-term scheduling determina partizione in cui caricare un processo, esistono due tecniche base:
\begin{itemize}
    \item una coda per ciascuna partizione:
        \begin{itemize}
            \item alla creazione un processo viene assegnato a coda relativa alla partizione più piccola in grado di contenerlo
            \item limitata flessibilità (possono esserci più processi in coda e partizioni non utilizzate)
        \end{itemize}
    \item una coda singola: processo allocato nella partizione libera più piccola tra quelle in grado di contenerlo (in caso non esistano $\rightarrow$ swap-out)
\end{itemize}
\includegraphics[width=\textwidth]{OS_images/PF_Malloc.png}
\subsubsection{Supporto per Rilocazione}
MMU consiste di registri base e limite: \\
\includegraphics[width=0.6\textwidth]{OS_images/Registri_MMU.png}
\subsubsection{Valutazione}
\begin{itemize}
    \item Vantaggi
        \begin{itemize}
            \item semplice da implementare
            \item basso carico di gestione da parte di S.O.
        \end{itemize}
    \item Svantaggi
        \begin{itemize}
            \item grado di multiprogrammazione fissato
            \item si genera \textbf{frammentazione}
                \begin{itemize}
                    \item \textit{interna} (se dimensione della partizione superiore a dimensione del processo che vi si carica)
                    \item \textit{esterna} (se vi sono partizioni non utilizzato perché non soddisfano esigenze dei processi in coda)
                \end{itemize}
        \end{itemize}
\end{itemize}
\subsection{Partizioni variabili}
Si vorrebbe partizionare memoria in parti di dimensioni corrispondenti a dimensioni dei processi che verranno caricati, ottenendo così eliminazione della frammentazione interna. Questo risultato si ottiene non fissando a priori numero e dimensioni delle partizioni.
\subsubsection{Allocazione della Memoria}
Si vede la memoria come un insieme di \textbf{buchi} (hole), ognuno di questi è una sezione di memoria. S.O. mantiene informazioni su partizioni assegnate e buchi disponibili. Un processo viene allocato in uno dei buchi in grado di contenerlo. \\
Dato che possono esserci più buchi quando un processo deve essere caricato c'è il problema di scegliere in che buco allocarlo, 3 soluzioni:
\begin{itemize}
    \item First-Fit (si occupa primo buco abbastanza grande che si trova)
    \item Best-Fit (si occupa buco più piccolo tra quelli in grado di contenre il processo)
    \item Worst-Fit (si occupa buco più grande)
\end{itemize}
\subsubsection{Supporto per Rilocazione}
MMU con stessa struttura della tecnica con partizioni fisse
\subsubsection{Valutazione}
\begin{itemize}
    \item Vantaggi
        \begin{itemize}
            \item non esiste frammentazione interna
            \item grado di multiprogrammazione non prefissato
        \end{itemize}
    \item Svantaggi
        \begin{itemize}
            \item frammentazione esterna può essere problematica
        \end{itemize}
\end{itemize}
\subsubsection{Compattazione}
Soluzione a problema della frammentazione esterna: rilocare processi e compattare soazio libero ottenendo unico buco (possibile solo con rilocazione dimanica) \\
\includegraphics[width=0.6\textwidth]{OS_images/Comp_VS_Frag_Mem.png}
\subsection{Buddy System}
Compromesso tra partizioni fisse e variabili. Sia $2^m$ dimensione del più piccolo buco che si vuole allocare e $2^M$ tutta la memoria, allora:
\begin{itemize}
    \item memoria vista come insieme di liste di blocchi di dimensione $2^k$ con $m \leq k \leq M$
    \item memoria disponibile è quindi insieme di blocchi di dimensione $2^i$ (per ogni i)
    \item all'inizio la lista dei blocchi di dimensione $2^M$ coniene un solo blocco (tutta la memoria) mentre altre liste sono vuote
    \item quando giunge richiesta di dimensione $s$, si crea blocco con dimensione sufficiente (potenza di 2):
        \begin{itemize}
            \item se $2^{M-1} < s \leq 2^M$ si alloca intera memoria
            \item altrimenti blocco di dimensione $2^M$ viene diviso in due blocchi di dimensione $2^{M-1}$ e si verifica se $2^{M-2} < s < 2^{M-1}$, ...
        \end{itemize}
    \item quando un processo rilascia suo blocco questo viene inserito nell'opportuna lista di blocchi liberi
    \item quando due blocchi adiacenti di dimensione $2^i$ si liberano, possono essere fusi in un blocco di dimensione $2^{i+1}$
\end{itemize}
\subsubsection{Valutazione}
\begin{itemize}
    \item Vantaggi: compattazione si effettua scorrendo lista dei blocchi
    \item presenza di frammentazione interna (non è detto che processi abbiano sempre dimensione $2^i$)
\end{itemize}
\includegraphics[width=0.6\textwidth]{OS_images/Buddy_System_Example.png}
\section{Paginazione}
\subsection{Generalità}
Introdotta per eliminare frammentazione esterna.
\begin{itemize}
    \item Idea: permettere che un processo \textbf{non sia allocato in modo contiguo}
    \item memoria fisica organizzata in blocchi di dimensione fissa (\textbf{frame})
    \item memoria logica organizzata in blocchi della stessa dimensione detti \textbf{pagine}
    \item S.O. tiene traccia dei frame liberi
    \item per caricare un processo che occupa $n$ pagine si devono usare $n$ frame liberi
    \item traduzione degli indirizzi basata su \textbf{page table} (una per processo)
    \item solo frammentazione interna e solo nell'ultima pagina di ogni processo
\end{itemize}
\subsection{Architettura}
Indirizzo logico viene diviso in due parti:
\begin{itemize}
    \item \textbf{Numero di pagina} $p$: indice nella page table per determinare numero del frame in memoria fisica
    \item \textbf{Offset (spiazzamento)} $d$: combinato con numero del frame determina indirizzo fisico
\end{itemize}
\includegraphics[width=0.6\textwidth]{OS_images/Paginazione.png} \\
Si tiene traccia di frame liberi / occupati con una \textbf{tabella dei frame}: \\
\includegraphics[width=0.6\textwidth]{OS_images/Frame_Table.png}
\subsubsection{Gestione dei processi}
Gestione processi si complica in quanto: \begin{itemize}
    \item \textbf{per ogni processo} va gestita una page table: nel PCB si inserisce un puntatore a page table
    \item context switch diventa più oneroso (vanno garantite informazioni sulla paginazione)
    \item swapping ifluenza contenuto della tabella dei frame
\end{itemize}
\subsection{Page Table}
Può essere realizzata in più modi:
\begin{itemize}
    \item insieme di registri dedicati
        \begin{itemize}
            \item efficiente ma non adattabile con tabelle grandi
            \item context switch più complesso
        \end{itemize}
    \item tabella in memoria
        \begin{itemize}
            \item localizza tramite \textbf{page-table base register (PTBR)}
            \item un solo registro da gestire in context switch
            \item difetto:per accedere ogni indirizzo servono \textit{due} accessi a memoria
        \end{itemize}
    \item utilizzo di cache speciale (\textbf{translationlook-aside buffer (TLB)}), memoria associativa key-value
\end{itemize}
\subsubsection{Page Table con TLB}
\begin{itemize}
    \item TLB contiene solo \textit{parte} della page table
    \item per tradurre un indirizzo si accede alla TLB
    \item in caso di \textbf{TLB-miss} si dovrà accedere alla vera tabella in memoria
    \item gestione di quale parte della tabella duplicare nella TLB è fatta con tecniche usate per gestire la cache (FIFO, LRU, ...)
    \item percentuale di volre che si verifica un \textbf{TLB-hit} è detta \textbf{hit ratio}
    \item durante context switch la TLB va svuotata
\end{itemize}
\includegraphics[width=0.6\textwidth]{OS_images/Page_Table_TLB.png}
\subsubsection{Protezione}
Ad ogni elemento della table si associa un \textbf{bit di validità} per individuare quali pagine sono nello spazio logico del processo.\\
Altri bit possono essere nella table per indicare che un particolare frame può essere oggetto di lettura, scrittura,... da parte di un processo.
\subsubsection{Struttura}
\begin{itemize}
    \item \textbf{Paginazione Gerarchica} $\rightarrow$ anche tabella delle pagine viene paginata. Esempio di organizzazione \textbf{a 2 livelli}: la parte "numero di pagina" $p$ dell'indirizzo logico è a sua volta divisa in due porzioni $p1$ e $p2$ \\
    \begin{minipage}{0.45\textwidth}
        \includegraphics[width=\textwidth]{OS_images/Gerarchia1.png}
    \end{minipage}
    \hfill
    \begin{minipage}{0.5\textwidth}
        \includegraphics[width=\textwidth]{OS_images/Gerarchia2.png}
    \end{minipage}
    Sono possibili organizzazioni con \textbf{più livelli} (architetture ARM usano paginazione gerarchica fino a 4 livelli) \\
    \textbf{Esempio}: considera indirizzi logici di 32 bit con pagine da 4K, l'indirizzo è suddiviso in:
        \begin{itemize}
            \item 20 bit per numero di pagina
            \item 12 bit di offset all'interno di una pagina ($d$)
        \end{itemize}
    Dato che anche page table è paginata (supponiamo in pagine di 1K), i 20 bit sono divisi in:
        \begin{itemize}
            \item 10 bit per outer-table ($p_1$)
            \item 10 bit per offset ($p_1$) all'interno della pagina della page table
        \end{itemize}
    \item \textbf{Page Table con Hashing} $\rightarrow$ in caso di aumento di livelli in paginazione gerarchica aumenta numero di accessi a memoria necessari per tradurre un indirizzo, una soluzione può essere l'uso di \textbf{organizzazione hash}:
    \includegraphics[width=0.6\textwidth]{OS_images/Page_Table_HASH.png}
    \item \textbf{Page Table Invertita} $\rightarrow$ in caso di dimensione della table troppo grande (anche rispetto a numero di frame fisici a disposizione)
        \begin{itemize}
            \item tabella che memorizza solamente informazioni su pagine realmente utilizzate da processi
            \item unica per tutto il sistema
            \item un'entry per ogni frame fisico ($i$), questa contiene indirizzo logico della pagina ($p$) caricata in quel frame e indicazione del processo possessore ($pid$)
        \end{itemize}
    \includegraphics[width=0.6\textwidth]{OS_images/Page_Table_Invertita.png} \\
    \textbf{Minor memoria occupata} per la table ma \textbf{maggior tempo di traduzione}: va cercata coppia $(pid, p)$ nella tabella ormai ordinata rispetto a indirizzi fisici ($i$) $\rightarrow$ rimedio: implementare inverted page table con hashing.
\end{itemize}
\subsubsection{Pagine Condivise}
Tecniche viste permettono di mappare pagine logiche in pagine fisiche ma anche di \textbf{mappare diverse pagine logiche nella stessa pagina fisica}:
\begin{itemize}
    \item si permette condivisione di dati
    \item si permette condivisione di codice (evitando ripetizioni di codice comune come editor, librerie dinamiche, compilatori...)
    \item codice condiviso deve essere \textbf{rientrante}: non in grado di auto-modificarsi
\end{itemize}
\includegraphics[width=0.6\textwidth]{OS_images/Pag_Cond_Example.png}
\section{Segmentazione}
\subsection{Generalità}
Tecniche finora viste sono accomunate da considerazione della memoria come unico vettore di locazioni
\begin{itemize}
    \item utente tende a imporre una struttura alla memoria: insieme di \textbf{segmenti}, ognuno con \textbf{nome}, \textbf{dimensione}, e \textbf{specifica funzione} (main program, codice procedura, stack, libreria...)
    \item indirizzi possono essere specificati fornendo nome del segmento e offset al suo interno
    \item riferimenti interni al segmento sono relativi all'inizio dello stesso
    \item \textbf{compilatore} compila programma dell'utente generando opportuni segmenti, altri possono essere introdotti dal \textbf{linker}, il \textbf{loader} fornisce "nomi" ai segmenti
\end{itemize}
\includegraphics[width=0.6\textwidth]{OS_images/Segmentazione.png} \\
Struttura logica del programma e sua allocazione
\subsection{Architettura}
\begin{itemize}
    \item Programma visto come insieme di segmenti (non per forza della stessa dimensione) %PEFFOH
    \item ogni segmento allocato in una porzione \textit{contigua} di memoria fisica
    \item traduzione degli indirizzi avviene tramite \textbf{segment table} che, per ogni segmento contiene:
        \begin{itemize}
            \item \textbf{base}: indirizzo fisico dell'inizio del segmento
            \item \textbf{limite}: lunghezza del segmento
        \end{itemize}
    \item la segment table è:
        \begin{itemize}
            \item nantenuta in registri locali
            \item o allocata in memoria e gestita traite \textbf{segment-table base register (STBR)} e il \textbf{segment-table length register (STLR)} %STeel Ball Run
        \end{itemize}
\end{itemize}
\includegraphics[width=0.6\textwidth]{OS_images/Arch_Segmentazione.png}
\subsection{Considerazioni}
\begin{itemize}
    \item \textbf{rilocazione è dinamica} (a livello di segmenti)
    \item si può prevedere \textbf{condivisione di segmenti} (es. librerie)
    \item allocazione presenta problema della \textbf{frammentazione esterna} (segmenti non hanno tutti stessa dimensione $\rightarrow$ stessi problemi di allocazione contigua)
    \item per allocare file si usano first-fit, best-fit...
    \item \textbf{protezione}: per ogni inserimento in segment table si usa un \textbf{bit di validità} e eventuali bit per privilegi / modalità d'accesso
    \item paginazione \textbf{trasparente all'utente} (necessità del S.O.) mentre segmentazione è \textbf{visibile} (richiesto che utente / compilatore descriva i segmenti)
\end{itemize}
\section{Segmentazione Paginata}
\subsection{Generalità}
Segmentazione e paginazione possono essere combinate per ottenere in un'unica organizzazione i pregi delle due tecniche:
\begin{itemize}
    \item spazio logico segmentato (riflette struttura logica del programma) come nella segmentazione
    \item ogni \textbf{segmento è paginato} (suddiviso in pagine della dimensione del frame fisico)
    \item per ogni processo si gestisce \textbf{una segment table} ed un certo numero di page table: \textbf{una per ogni segmento}
    \item si evita frammentazsione esterna a prezzo di una frammentazione (minore) interna (ultima pagina di ogni segmento)
\end{itemize}
\subsection{Architettura}
\includegraphics[width=0.6\textwidth]{OS_images/Arch_Seg_Paginata.png}
\end{document}
